\section*{Chapitre \ref{ch:rs:the-research-process} --- Le processus de recherche}

\begin{solution}{exo:rs:the-research-process:1}
$\mathrm{PPV} = 0.1 \times 0.8/(0.1 \times 0.8 + 0.9 \times 0.05) = 0.08/0.125 = 0.64$. Avec un a priori de 2~\%, $0.016/(0.016 + 0.049) = 0.246$~:
trois passages sur quatre sont faux.
\end{solution}

\begin{solution}{exo:rs:the-research-process:2}
$1 - \Phi(2)^{20} = 1 - 0.97725^{20} = 0.369$~; en moyenne $20 \times 0.02275 = 0.455$ passent.
\end{solution}

\begin{solution}{exo:rs:the-research-process:3}
$\Phi\bigl((-0.5 - 1.2)\sqrt{0.5}\bigr) = \Phi(-1.20) = 0.115$~: environ un semestre sur neuf serait aussi mauvais. Pas décisif à lui seul.
\end{solution}

\begin{solution}{exo:rs:the-research-process:4}
$Y = ((1.645 + 1.282)/0.8)^2 = 13.4$ ans, contre 2.7 ans pour 1.5 avec une puissance de 80~\%~: la durée croît comme l'inverse du carré
du ratio de Sharpe.
\end{solution}

\begin{solution}{exo:rs:the-research-process:5}
$30 + \binom{30}{2} = 30 + 435 = 465$ tests. La correction de Bonferroni exige $p \le 0.05/465 = 1.08 \times 10^{-4}$, soit $t \ge 3.70$.
\end{solution}

\begin{solution}{exo:rs:the-research-process:6}
Six idées réelles et 144 fausses. Après le premier \omterm{def:rs:the-research-process:gate}{jalon de validation}, 3.6 réelles et 7.2 fausses (PPV 0.333)~; après le deuxième, 2.16 et 0.36 (0.857)~;
après le troisième, 1.30 et 0.018 (0.986). Environ 1.3 stratégies réelles et une fausse tous les 55 ans parviennent en production.
\end{solution}

\begin{solution}{exo:rs:the-research-process:7}
\texttt{p\_best(2, 1000, 2, rho)} donne 0.420, 0.167 et 0.030 pour $\rho = 0.3$, 0.6 et 0.9. Vingt variantes indépendantes donnent 0.046~;
mille variantes équicorrélées à $\rho = 0.85$ donnent le même résultat. Même à 0.9, mille variantes valent environ treize
indépendantes~: la corrélation réduit une recherche, mais rarement à un seul essai.
\end{solution}

\begin{solution}{exo:rs:the-research-process:8}
L'\omterm{def:rs:the-research-process:prereg}{échantillon réservé} a cessé de l'être dès qu'il a servi à l'ajustement~: le 2.1 est un chiffre dans l'échantillon, et la seule preuve indépendante
est le 1.4 mesuré avant l'ajustement, lui-même un regard sur un \omterm{def:rs:the-research-process:prereg}{échantillon réservé} qui a ensuite été réutilisé. Deux paramètres ajustés sur trois ans sont aussi
essais que le journal doit compter. Le statut honnête est celui d'un seul \omterm{def:rs:the-research-process:gate}{jalon de validation} exploratoire franchi~; une nouvelle période, non touchée (ou du trading fantôme (shadow trading)
en direct) est nécessaire avant que l'affirmation «~deux \omterm{def:rs:the-research-process:gate}{jalons de validation} indépendants~» soit vraie.
\end{solution}

\begin{solution}{pb:rs:the-research-process:1}
\textbf{1.} Environ $\mathcal N(0, 1/2)$~: écart-type 0.707. \textbf{2.} $1 - \Phi(2\sqrt2) = 1 - \Phi(2.83) = 0.00234$.
\textbf{3.} $1\,000 \times 0.00234 = 2.34$. \textbf{4.} La question porte sur la probabilité qu'au moins une impressionne (et laquelle est
montrée)~; la corrélation modifie cette probabilité, non l'espérance. \textbf{5.} $1 - (1 - 0.00234)^{20} = 0.0458$.
\textbf{6.} Par la \cref{prop:rs:the-research-process:best} avec $\rho = 0.6$~: 0.0262. \textbf{7.} $1 - (1 - 0.0262)^{50} = 0.735$.
\textbf{8.} $1 - \Phi(2.83)^{1\,000} = 0.904$. \textbf{9.} 2.30, au-dessus du seuil «~impressionnant~».
\textbf{10.} $\Phi\bigl(2/\sqrt{252} \times \sqrt{503}\bigr) = 0.998$. \textbf{11.} Avec l'indice de référence, le meilleur attendu parmi vingt ratios nuls
ratios~: 0.822. \textbf{12.} Avec mille essais~: 0.336, soit un pile ou face pour que la stratégie batte le meilleur résultat de la recherche.
\textbf{13.} Le décompte doit couvrir tout essai qui aurait pu être présenté~; la corrélation réduit le nombre effectif, et le
ratio déflaté calculé avec le décompte brut est alors prudent, mais le nombre effectif exige les corrélations des variantes, ce que seul
le journal peut fournir. \textbf{14.} $2\sqrt Y \ge z_{0.99999} = 4.265$, donc $Y \ge 4.55$ ans.
\textbf{15.} $1 - \Phi(1) = 0.159$. \textbf{16.} Puissance $\Phi(1.5 - 1) = 0.691$~; $\mathrm{PPV} = 0.05 \times 0.691/(0.05 \times 0.691 + 0.95
\times 0.159) = 0.187$~: un \omterm{def:rs:the-research-process:prereg}{échantillon réservé} d'un an à ce seuil constitue un jalon faible. \textbf{17.} $1 - (1 - 0.159)^4 = 0.499$~: un tirage à
pile ou face. \textbf{18.} \emph{Résultat nommé~:} une recherche sans valeur de vingt variantes corrélées produit chaque semaine au moins un backtest
avec un ratio de Sharpe de 2 dans l'année avec une probabilité de 0.735, et le ratio de Sharpe déflaté de ce backtest, calculé avec les
mille essais journalisés, vaut 0.336. \textbf{19.} L'hypothèse et la date de son enregistrement~; le \omterm{def:rs:the-research-process:log}{nombre d'essais} de la famille et la
distribution des résultats de tous les essais~; les preuves au-delà de l'ajustement dans l'échantillon (données ultérieures, prédictions annexes, un \omterm{def:rs:the-research-process:prereg}{échantillon réservé}).
\textbf{20.} Il consigne la recherche, et la recherche est le plus grand risque non modélisé d'un backtest.
\end{solution}

\begin{solution}{iq:rs:the-research-process:1}
Combien de variantes ont été essayées, et à quel point étaient-elles corrélées (le \omterm{def:rs:the-research-process:log}{nombre d'essais}, donc le ratio déflaté)~? Quelle est la
\omterm{def:rs:the-research-process:hypothesis}{justification économique}, et a-t-elle été rédigée avant le test~? Quelles preuves existent au-delà de l'ajustement dans l'échantillon (données ultérieures, autres marchés, prédictions
annexes), et les frais, le coût d'emprunt et la capacité sont-ils réalistes~?
\iqlookfor{la recherche avant les statistiques~; une justification assortie de prédictions annexes testables~; les coûts.}
\end{solution}

\begin{solution}{iq:rs:the-research-process:2}
D'après Bayes, la part de stratégies réelles parmi celles qui passent est $\pi(1 - \beta)/(\pi(1 - \beta) + (1 - \pi)\alpha)$. Avec un a priori de 5~\%,
un niveau de 5~\% et une puissance de 50~\%, elle vaut 0.345~: la plupart des passages sont faux quel que soit le niveau, à moins que l'a priori ne monte ou que des jalons indépendants
soient ajoutés. Le niveau contrôle les faux positifs par idée fausse testée, et non la proportion de stratégies fausses parmi les stratégies retenues.
\iqlookfor{la valeur prédictive positive, et le fait que les a priori et les jalons dominent.}
\end{solution}

\begin{solution}{iq:rs:the-research-process:3}
Tout usage de l'échantillon qui influence un choix~: l'ajuster, le consulter puis revenir modifier le modèle, le relancer après un
résultat décevant, ou choisir la période de l'échantillon en connaissance de cause. Chaque regard augmente la taille effective de l'\omterm{def:rs:the-research-process:prereg}{échantillon réservé} (cinq
regards à 5~\% donnent 22.6~\%).
\iqlookfor{qu'un \omterm{def:rs:the-research-process:prereg}{échantillon réservé} est consommé par l'influence, et non par une simple lecture.}
\end{solution}

\begin{solution}{iq:rs:the-research-process:4}
Pas sur cette seule preuve~: si le ratio réel valait 1.5, une année à zéro ou en dessous a une probabilité $\Phi(-1.5) = 0.067$. Examiner ce qui
est mesuré plus vite que les rendements~: les exécutions, les coûts et la rotation sont-ils ceux que prévoyait le simulateur~? Le coefficient d'information du signal
en direct est-il conforme au backtest~? Appliquer le \omterm{def:rs:the-research-process:gate}{critère d'abandon} fixé avant la mise en production~; si aucun n'a été fixé, en fixer un maintenant et
inclure une durée minimale.
\iqlookfor{le bruit d'un ratio de Sharpe sur un an, et les diagnostics au-delà du P\&L.}
\end{solution}

\begin{solution}{iq:rs:the-research-process:5}
Par entrée~: la famille, le type (hypothèse ou essai), l'horodatage fourni par l'appelant, les paramètres, les métriques, la version du code, l'instantané de données et
l'empreinte de l'entrée précédente, l'empreinte propre de l'entrée étant un hachage de l'ensemble de ces éléments. Stockage en ajout seul~; la vérification parcourt
la chaîne. Une hypothèse par famille, enregistrée avant tout essai~; les essais qui la précèdent signalent la famille a posteriori.
\iqlookfor{chaînage de hachages, horodatages reproductibles, et comptage des essais par famille d'hypothèses.}
\end{solution}

\begin{solution}{iq:rs:the-research-process:6}
Les vendeurs ou acheteurs qui exigent l'immédiateté éloignent les cours de la valeur, et les fournisseurs de liquidité qui absorbent le flux sont rémunérés par le
retournement ultérieur. Prédictions~: le retournement est plus fort pour les actions moins liquides~; plus fort lorsque la liquidité de marché est rare
(forte volatilité, bilans des dealers contraints)~; plus faible, voire nul, lorsque le mouvement porte une information (jours de résultats), et
s'atténuant à mesure que le coût de l'immédiateté pour le marché baisse au fil des années.
\iqlookfor{une justification qui dit qui paie, et des prédictions annexes capables de la réfuter.}
\end{solution}
